\section{Описание}
Требуется реализовать программную библиотеку на основе красно-чёрного дерева и использовать её для создания программы-словаря. В словаре каждому ключу, представляющему собой регистронезависимую последовательность букв английского алфавита длиной не более $256$ символов, ставится в соответствие число типа \texttt{unsigned long long}.

Поддерживаются операции добавления, удаления, поиска, сохранения и загрузки словаря из бинарного файла. При этом программа должна обрабатывать входной поток до конца файла, а при возникновении системных ошибок выводить сообщение, начинающееся с \enquote{ERROR:}.

В качестве структуры данных было выбрано красно-чёрное дерево, так как оно обеспечивает логарифмическое время выполнения основных операций: поиска, вставки и удаления. Для обеспечения регистронезависимости все ключи перед обработкой приводятся к верхнему регистру.

Особое внимание было уделено корректной сериализации и десериализации дерева. Сохранение выполняется в компактном бинарном формате с использованием маркеров для пустых и непустых узлов, а загрузка выполняется с проверкой корректности формата файла. Если при загрузке возникает ошибка, рабочее дерево остаётся без изменений.

\pagebreak

\section{Исходный код}
Программа состоит из трёх основных частей: вспомогательных структур для результата выполнения команд, класса красно-чёрного дерева и разбора команд из входного потока.

Для хранения результата операций добавления, удаления и поиска была введена структура \texttt{ReturnData}. Она хранит статус операции и, при необходимости, значение, связанное с найденным ключом.

Класс \texttt{RBTree} содержит вложенную структуру \texttt{Node}, в которой хранятся ключ, значение, цвет узла и ссылки на потомков и родителя. Для удобства реализованы вспомогательные методы, позволяющие определять положение узла относительно родителя, находить деда, брата и дядю, а также выполнять доступ к ссылке родителя на текущий узел.

Основные повороты дерева реализованы отдельными методами \texttt{leftRotation} и \texttt{rightRotation}. На их основе выполняется балансировка после вставки и удаления. При вставке используется стандартный алгоритм восстановления свойств красно-чёрного дерева через перекраску и повороты. При удалении учитываются все возможные случаи: удаление листа, узла с одним потомком и узла с двумя потомками, когда удаляемый элемент заменяется минимальным элементом правого поддерева.

Для сохранения дерева используется обход в прямом порядке: для каждого узла в файл записываются маркер узла, цвет, длина ключа, сам ключ и значение. Пустые поддеревья кодируются специальным маркером. При загрузке дерево восстанавливается рекурсивно по тем же маркерам. Если файл содержит некорректные данные или в нём присутствуют лишние байты, загрузка завершается ошибкой.

Ниже приведены основные объявления структур и класса.

\begin{lstlisting}[language=C++]
#include <fstream>
#include <iostream>
#include <string>

enum Side {
    left,
    right
};

struct ReturnData {
    enum Status {
        ok,
        exist,
        noSuchWord
    } status;

    unsigned long long value;
};

class RBTree {
    struct Node {
        enum Color {
            red,
            black
        };

        std::string key;
        unsigned long long value;
        Color color;
        Node* left = nullptr;
        Node* right = nullptr;
        Node* parent = nullptr;

        Node(const std::string& key, unsigned long long value,
             Node* parent, Color color);
        Node(std::ifstream& file);

        void update(const std::string& key, unsigned long long value);
        bool isLeftSon();
        Node* grandpa();
        Node* bro();
        Node* uncle();
        Node*& parentLink();
        void save(std::ofstream& file) const;
    };

    Node* root = nullptr;

    void destroy(Node* node);
    Node* leftRotation(Node* node);
    Node* rightRotation(Node* node);
    void insertBalance(Node* node);
    ReturnData::Status _insert(Node* node,
                               const std::string& key,
                               unsigned long long value);
    void removeBalance(Node* parent, Side side);
    Node* minNode(Node* node);
    void _remove(Node* node);
    Node* _find(Node* node, const std::string& key);
    static void toUpper(std::string& str);
    void _save(Node* node, std::ofstream& file);
    Node* _load(std::ifstream& file);

public:
    RBTree();
    ~RBTree();

    ReturnData::Status insert(std::string key, unsigned long long value);
    ReturnData::Status remove(std::string key);
    ReturnData find(std::string key);
    ReturnData::Status save(const std::string& filename);
    ReturnData::Status load(const std::string& filename);
};

struct Command {
    enum Function {
        insert,
        remove,
        find,
        save,
        load
    };

    Function func;
    std::string key;
    unsigned long long value;
};
\end{lstlisting}

Для разбора входных команд используется структура \texttt{Command}. Она определяет тип операции по первому токену строки и, в зависимости от него, считывает ключ и значение. Поддерживаются команды добавления, удаления, поиска, сохранения и загрузки.

Главная функция программы организует цикл чтения команд до конца входного потока. Для каждой операции вызывается соответствующий метод класса \texttt{RBTree}. Все исключения перехватываются, после чего выводится строка с текстом ошибки.

\begin{longtable}{|p{6.8cm}|p{8.2cm}|}
\hline
\rowcolor{lightgray}
\multicolumn{2}{|c|}{main.cpp}\\
\hline
\texttt{struct ReturnData}&Структура для хранения результата выполнения операций.\\
\hline
\texttt{class RBTree}&Реализация красно-чёрного дерева, поддерживающего вставку, удаление, поиск, сохранение и загрузку.\\
\hline
\texttt{struct Command}&Разбор строки входного файла и определение типа команды.\\
\hline
\texttt{void toUpper(std::string\& str)}&Преобразование ключа к верхнему регистру для обеспечения регистронезависимости.\\
\hline
\texttt{Node* leftRotation(Node* node)}&Левый поворот поддерева.\\
\hline
\texttt{Node* rightRotation(Node* node)}&Правый поворот поддерева.\\
\hline
\texttt{void insertBalance(Node* node)}&Восстановление свойств красно-чёрного дерева после вставки.\\
\hline
\texttt{void removeBalance(Node* parent, Side side)}&Восстановление свойств красно-чёрного дерева после удаления.\\
\hline
\texttt{ReturnData::Status insert(...)}&Добавление нового слова в словарь.\\
\hline
\texttt{ReturnData::Status remove(...)}&Удаление слова из словаря.\\
\hline
\texttt{ReturnData find(...)}&Поиск слова и вывод его значения.\\
\hline
\texttt{ReturnData::Status save(...)}&Сохранение дерева в бинарный файл.\\
\hline
\texttt{ReturnData::Status load(...)}&Загрузка дерева из бинарного файла с заменой текущего словаря при успехе.\\
\hline
\end{longtable}

\pagebreak

\section{Консоль}
\begin{alltt}
root@66a8fcab42ba:/workspaces/discran-labs/laba-2# cat test.txt
+ a 1
+ A 2
+ aaaaaaaaaaaaaaaaaaaaaaaaaaaaaaaaaaaaaaaaaaaaaaaaaaaaaaaaaaaaaaaaaaaaaaaaaaaaaaaaaaaaaaaaaaaaaaaaaaaaaaaaaaaaaaaaaaaaaaaaaaaaaaaaaaaaaaaaaaaaaaaaaaaaaaaaaaaaaaaaaaaaaaaaaaaaaaaaaaaaaaaaaaaaaaaaaaaaaaaaaaaaaaaaaaaaaaaaaaaaaaaaaaaaaaaaaaaaaaaaaaaaaaaaaaaaaaaa 18446744073709551615
aaaaaaaaaaaaaaaaaaaaaaaaaaaaaaaaaaaaaaaaaaaaaaaaaaaaaaaaaaaaaaaaaaaaaaaaaaaaaaaaaaaaaaaaaaaaaaaaaaaaaaaaaaaaaaaaaaaaaaaaaaaaaaaaaaaaaaaaaaaaaaaaaaaaaaaaaaaaaaaaaaaaaaaaaaaaaaaaaaaaaaaaaaaaaaaaaaaaaaaaaaaaaaaaaaaaaaaaaaaaaaaaaaaaaaaaaaaaaaaaaaaaaaaaaaaaaaaa
A
- A
a
root@66a8fcab42ba:/workspaces/discran-labs/laba-2# g++ main.cpp -Wall
root@66a8fcab42ba:/workspaces/discran-labs/laba-2# ./a.out < test.txt
OK
Exist
OK
OK: 18446744073709551615
OK: 1
OK
NoSuchWord
root@66a8fcab42ba:/workspaces/discran-labs/laba-2# 
\end{alltt}

В данном примере выполняются операции добавления, поиска, повторного добавления уже существующего слова. Программа корректно обрабатывает команды до конца входного файла и выводит результат каждой операции в соответствии с условием лабораторной работы.
\pagebreak